%! Author = sriadityavedantam
%! Date = 4/8/26

\section{Chapter 9}

\begin{problem}[Exercise 1]{
    Given $G$ is a group of order $25$, show that $G$ is cyclic or an element in $G$ has order $5$.
}
    Given $G$ is a finite group, by Corollary 8.6 the order of every $a \in G$ must divide the order of $G$.
    The possible factors of $25$ are $1,5,25$.
    If $\ord(a) = 1$, then $a = e$.
    If $\ord(a) = 5$, then $G$ has an element of order $5$.
    If $\ord(a) = 25$, then there is one generating element of $G$, making it cyclic.
    Thus this proves the claim.
\end{problem}

\begin{lemma}[Reduction of Possible Divisors of $2n$]{
    The maximum divisor of $2n$ is $n$.
}
    Note that the total cap to integer divisors is $\lfloor \sqrt{2n} \rfloor$.
    What can we do to reduce this range in relation to $n$.
    Could we have a possible divisor of $2n$ put into relation with $n$.
    That is the question this lemma attempts to answer.
    By the Fundamental Theorem of Arithmetic, all divisors of $n$ divide $2n$.
    However, if $n$ is built up of primes, $n = p_1 p_2 \dots p_k$, then the only difference between $n$ and $2n$ is the prime $2$.
    Thus $2n = 2p_1 p_2 \dots p_k$.
\end{lemma}

\begin{problem}[Exercise 2]{
    Given $G$ is an abelian group of order $2n$, with $n$ odd, suppose $G$ contains more than one element of order $2$.
}
    By Exercise 8.1.31, since $\ord(G) = 2n$, which is even, $G$ contains at least one element of order $2$.
    By the previous lemma, there are no divisors of $2n$ greater than $n$ that are relevant here.
    Thus the only divisors of $n$ have to be odd by the Fundamental Theorem of Arithmetic, leaving only one even divisor, namely $2$.
\end{problem}

\begin{problem}[Exercise 3]{
    Prove the following.
}
    \textbf{(a).}
    Suppose $a$ and $b$ have order $3$ in a group and $a^2 = b^2$.
    Then
    \begin{align*}
        a^2 a^2 &= b^2 b^2\\
        a^4 &= b^4\\
        a^3 a &= b^3 b\\
        ea &= eb\\
        a &= b.
    \end{align*}
    Thus by these statements, $a = b$.

    \textbf{(b).}
    Given $G$ is a finite group, and for $a \in G$, $\ord(a) = 3$, then
    \begin{align*}
    (a^2)^3 &= (a^3)^2\\
        &= e^2\\
        &= e.
    \end{align*}
    As we proved in part \textbf{(a)}, $a$ has order $3$ already.
    Thus since $a$ and $a^2$ have order $3$, all elements of order $3$ occur in pairs, being included in the set
    \[
        \{b,b^2 : \forall b \in G \text{ and } \ord(b) = 3\}.
    \]
    Hence the number of elements of order $3$ is even.
\end{problem}

\begin{problem}[Exercise 4]{
    Given $p$ prime and $p \mid \ord(G)$, let $S \coloneqq \{a \in G : \ord(a) = p\}$.
}
    By Exercise 8.1.34, there are exactly $p-1$ distinct elements of order $p$ in the subgroup of $G$ generated by $\langle a \rangle$.
    Define a relation on $S$ such that $\langle a \rangle = \langle b \rangle$, denoted $a \sim b$.
    This relation creates clearly\footnote{Allowed to state by Dr. Chastkofsky.} distinct classes with each class containing $p-1$ elements,
    \[
        [a] = \{b \in S : a \sim b\}.
    \]
    Hence $(p-1) \mid |S|$.
    In fact, we learn that the number of these elements is $n(p-1)$, with $n$ counting the number of distinct classes under this relation.
\end{problem}

\begin{problem}[Exercise 5]{
    Denote $Z = Z(G)$.
    Show that if $C \leq Z$, then $C \trianglelefteq G$.
}
    Given $G$ is a group, suppose we have an element $c \in Z$.
    Then by the definition of $Z$, we have that $c$ commutes with all elements in $G$.
    Thus we have that $Z$ is a normal subgroup of $G$.
    Suppose $C \leq Z$.
    Then every element in $C$ still commutes with every element in $G$.
    Thus subgroup $C$ is normal in $G$.
\end{problem}

\begin{problem}[Exercise 6]{
    Show that the center $Z(G)$ is a characteristic subgroup.
}
    Given group $G$, we have previously proved that $Z$ is a normal subgroup of $G$.
    Let $f$ be an automorphism of $G$.
    Given an element $c \in Z$, then
    \begin{align*}
        f(c)f(a) &= f(ca)\\
        &= f(ac)\\
        &= f(a)f(c).
    \end{align*}
    Since $f$ is an automorphism, this function is an isomorphism, thus unique and has an inverse.
    Thus $f(c) \in Z$, since $f(c) \in G$ and $f(a) \in G$ for all $a \in G$.
    Thus $Z$ is a characteristic subgroup.
\end{problem}

\begin{problem}[Exercise 7]{
    Given $K$ is a normal subgroup of order $2$ in group $G$, show that $K \subseteq Z(G)$.
}
    Given $k \in G$ with $\ord(k) = 2$, let $K = \{e,k\}$.
    Suppose $K \not\subseteq Z$.
    Then $k \in K$ is not commutative with all elements in $G$.
    We know since $\ord(k) = 2$, that $k$ commutes with itself.
    By Theorem 8.11, $Ka = aK$, thus by weak commutativity $ka = ak_1$ for all $a \in G$.
    However, by our assumption this is only possible if $k$ is not central, which gives a contradiction with the fact that $K$ has only two elements and is normal.
    Thus $K \subseteq Z$.
\end{problem}

\begin{problem}[Exercise 8]{
    Let $N$ and $K$ be subgroups of a group $G$.
}
    \textbf{(a).}
    Given $N$ and $K$ are subgroups of group $G$ and $N$ is normal in $G$, suppose $NK = \{nk : n \in N, k \in K\}$.
    Let $a = n_1k_1$ and $b^{-1} = k_2^{-1}n_2^{-1}$.
    Since $k_1k_2^{-1} \in K$ due to closure, then $k_1k_2^{-1}n_2^{-1} = n_3k_1k_2^{-1}$ by weak commutativity.
    Thus
    \begin{align*}
        ab^{-1} &= n_1k_1k_2^{-1}n_2^{-1}\\
        &= n_1n_3k_1k_2^{-1}.
    \end{align*}
    Since $n_1n_3 \in N$ and $k_1k_2^{-1} \in K$, then $ab^{-1} \in NK$.

    \textbf{(b).}
    Given $N$ and $K$ are normal subgroups of $G$, then for $a = n_1k_1$ and $b^{-1} = k_2^{-1}n_2^{-1}$, we have
    \begin{align*}
        ab^{-1} &= n_1k_1k_2^{-1}n_2^{-1}\\
        &= n_1k_1n_2^{-1}n_2k_2^{-1}n_2^{-1}\\
        &= n_1k_1n_2^{-1}(n_2k_2^{-1}n_2^{-1})\\
        &= n_1k_1n_2^{-1}k_3\\
        &= n_1n_2^{-1}n_2k_1n_2^{-1}k_3\\
        &= n_1n_2^{-1}(n_2k_1n_2^{-1})k_3\\
        &= n_1n_2^{-1}k_4k_3.
    \end{align*}
    Since $n_1n_2^{-1} \in N$ and $k_4k_3 \in K$, then $ab^{-1} \in NK$.
\end{problem}

\begin{problem}[Exercise 9]{
    Given $K$ and $N$ are normal subgroups of group $G$ such that $K \cap N = \langle e \rangle$, show that $kn = nk$ for all $k \in K$ and $n \in N$.
}
    Suppose we have $a = k^{-1}n^{-1}kn$.
    By using properties of normal subgroups, we have that $a \in N$ but also that $a \in K$, since
    \[
        (k^{-1}n^{-1}k)n = k^{-1}(n^{-1}kn).
    \]
    By the end of it, $a = e$, which satisfies $K \cap N$.
    Then
    \begin{align*}
        k^{-1}n^{-1}kn &= e\\
        (nk)^{-1}kn &= e\\
        nk(nk)^{-1}kn &= nke\\
        kn &= nk.
    \end{align*}
    Thus $nk = kn$.
\end{problem}

\begin{problem}[Exercise 10]{
    Let $N$ be a subgroup of index $2$ in $G$.
}
    \textbf{(a).}
    Given $N$ is a subgroup of $G$ of index $2$, then by Theorem 8.4, there are only two distinct cosets.
    Thus for $a \notin N$, the set $Na$ is exactly the set of all elements not in $N$.

    \textbf{(b).}
    Suppose for $a \notin N$ and $n \in N$, then $a^{-1}Na$ is a subgroup with the same order as $N$.
    Since $N$ has index $2$, there is only one subgroup of that index containing the identity.
    Thus it must be the case that $a^{-1}Na \subseteq N$.
    Therefore $N$ is a normal subgroup by Theorem 8.11.
\end{problem}

\begin{problem}[Exercise 11]{
    Given $\ord(H) = n$ and $H$ is the only subgroup of this order, show that $H$ is normal.
}
    By Exercise 7.4.20, $a^{-1}Ha \leq G$ and $H \cong a^{-1}Ha$, given that $a \in G$.
    Thus the orders of $H$ and $a^{-1}Ha$ are the same due to this isomorphism.
    But we are given that $H$ is the only subgroup of order $n$.
    Thus $a^{-1}Ha = H$.
    Hence $H$ is normal.
\end{problem}

\begin{problem}[Exercise 12]{
    Let $T$ be a subgroup of $G$ and suppose $Ta \in G/T$.
}
    By properties of right cosets,
    \begin{align*}
    (Ta)^n &= T(a^n).
    \end{align*}
    If $a \in T$ and $\ord(a) = n$, then $a^n = e$, thus
    \[
        Te = T.
    \]
    If $\ord(a) = \infty$, then $T(a^n)$ will never have finite order for any $n \in \n$.
\end{problem}

\begin{problem}[Exercise 13]{
    Given that $G$ is a simple group and $f : G \to H$ is a surjective homomorphism of groups, determine the possibilities for $H$.
}
    Let $K$ be the kernel of this function.
    By Theorem 8.4.16, $K$ is a normal subgroup of $G$.
    However, since $G$ is simple, this could either mean the trivial group or $G$ itself.
    If $K$ is trivial, then $K = \langle e \rangle$, thus by Theorem 8.4.17, $f$ is an isomorphism due to injectivity.
    If $K$ equals $G$ itself, then all elements in $G$ map to $e_H$, which means $H = \langle e \rangle$.
\end{problem}

\begin{lemma}[Exercise 7.2.33]{
    Assume that $a,b \in G$ and $ab = ba$.
    If $\ord(a)$ and $\ord(b)$ are relatively prime, prove that $ab$ has order $\ord(a)\ord(b)$.
}
    Suppose $\ord(a) = n$ and $\ord(b) = m$.
    If $m,n$ are relatively prime, then by Theorem 7.9, as long as $n \mid k$ in $a^k$, and $m \mid \ell$ in $b^\ell$, then $n$ and $m$ divide $\lcm(m,n) = nm$.
    Thus $\ord(ab) = \ord(a)\ord(b) = mn$.
\end{lemma}

\begin{corollary}[Corollary of Exercise 7.2.33]{
    The order of the product of two elements is $\ord(ab) = \lcm[\ord(a),\ord(b)]$.
}
    This follows from the previous lemma in the relatively prime case.
\end{corollary}

\begin{problem}[Exercise 14]{
    Let
    \[
        K = \{a \in G : \ord(a) \leq 2\}
        \qquad\text{and}\qquad
        H = \{a^2 : a \in G\}.
    \]
}
    \textbf{(a).}
    Note that the only element of order $1$ is the identity, thus $K$ is closed under the identity.
    Given how $\ord(a) = 2$, it is also closed under inverses due to the inverse also having the same order as the element.
    Suppose $a,b \in K$.
    Then $ab^{-1}$ has order less than or equal to $2$ due to $\lcm(\ord(a),\ord(b)) \leq 2$ by Corollary 9.1.1.
    Thus $K$ is closed under multiplication, hence a subgroup of $G$.

    \textbf{(b).}
    $H$ is closed under identity since $e^2 = e \in H$.
    $H$ is closed under inverses since given $a \in G$, $a^2,(a^{-1})^2 \in H$, but $(a^{-1})^2 = (a^2)^{-1}$.
    Given $a,b \in G$, then
    \[
        (ab^{-1})^2 = (b^{-1})^2a^2 = a^2(b^{-1})^2 \in H.
    \]
    Thus $H$ is closed under multiplication.
    Hence $H$ is a subgroup of $G$.

    \textbf{(c).}
    Define a map $f : G \to H$ by $x \mapsto x^2$ with kernel $K$.
    We claim $f$ is a homomorphism.
    Given $x,y \in G$,
    \begin{align*}
        f(xy^{-1}) &= (xy^{-1})^2\\
        &= (y^{-1})^2x^2\\
        &= x^2(y^{-1})^2\\
        &= f(x)f(y^{-1}).
    \end{align*}
    Since $H$ is composed of all squares of $G$, then $f$ is surjective.
    Given that
    \[
        \ker(f) = \{x \in G : f(x) = e_H\},
    \]
    note that $f(x) = x^2$, which follows our requirement for $K$ being the kernel holding all orders of $2$ or less.
    Hence by the First Isomorphism Theorem, $G/K \cong H$.
\end{problem}

\begin{problem}[Exercise 15]{
    Show that
    \[
        (\mathbb{Z} \times \mathbb{Z})/\langle(1,1)\rangle \cong \mathbb{Z}.
    \]
}
    Define $G = \mathbb{Z} \times \mathbb{Z}$ and define a map $f : G \to \mathbb{Z}$ by $(a,b) \mapsto a-b$.
    We claim that $f$ is a homomorphism.
    Given $(a,b),(c,d) \in G$,
    \begin{align*}
        f((a,b) + (c,d)) &= f(a+c,b+d)\\
        &= a+c-(b+d)\\
        &= a-b+c-d\\
        &= f(a,b) + f(c,d).
    \end{align*}
    Thus $f$ is a homomorphism.
    Suppose $(a,0) \in G$, then $f(a,0) = a-0 = a$, thus $f$ is surjective.
    Given
    \[
        \ker(f) = \{(a,b) \in G : f(a,b) = e_H\},
    \]
    this is only true when $a=b$, thus we have a kernel of $K = \langle(1,1)\rangle$.
    Hence by the First Isomorphism Theorem,
    \[
        (\mathbb{Z} \times \mathbb{Z})/\langle(1,1)\rangle \cong \mathbb{Z}.
    \]
\end{problem}

\begin{problem}[Exercise 16]{
    Show that
    \[
        GL(2,\rr)/SL(2,\rr) \cong \rr^*.
    \]
}
    Define a map
    \[
        f : GL(2,\rr) \to \rr^*
    \]
    by
    \[
        \begin{pmatrix}
            a & b\\
            c & d
        \end{pmatrix}
        \mapsto ad-bc.
    \]
    We claim $f$ is a homomorphism.
    Given
    \[
        \begin{pmatrix}
            a & b\\
            c & d
        \end{pmatrix},
        \begin{pmatrix}
            e & f\\
            g & h
        \end{pmatrix}
        \in GL(2,\rr),
    \]
    we have
    \begin{align*}
        f\left(
             \begin{pmatrix}
                 a & b\\
                 c & d
             \end{pmatrix}
             \begin{pmatrix}
                 e & f\\
                 g & h
             \end{pmatrix}
        \right)
        &=
        f\left(
             \begin{pmatrix}
                 ae+bg & af+bh\\
                 ce+dg & cf+dh
             \end{pmatrix}
        \right)\\
        &=
        (ae+bg)(cf+dh) - (af+bh)(ce+dg)\\
        &=
        (ad-bc)(eh-fg)\\
        &=
        f\left(
             \begin{pmatrix}
                 a & b\\
                 c & d
             \end{pmatrix}
        \right)
        f\left(
             \begin{pmatrix}
                 e & f\\
                 g & h
             \end{pmatrix}
        \right).
    \end{align*}
    Thus $f$ is a homomorphism.
    We claim that $f$ is surjective.
    Given
    \[
        \begin{pmatrix}
            x & 0\\
            0 & 1
        \end{pmatrix}
        \in GL(2,\rr),
    \]
    then
    \[
        f\left(
             \begin{pmatrix}
                 x & 0\\
                 0 & 1
             \end{pmatrix}
        \right)=x,
    \]
    thus $f$ is surjective.
    We have
    \[
        \ker(f)=\{A \in GL(2,\rr) : \det(A)=1\}=SL(2,\rr).
    \]
    Thus
    \[
        GL(2,\rr)/SL(2,\rr)\cong \rr^*.
    \]
\end{problem}

\begin{problem}[Exercise 17]{
    Given $K$ and $N$ both subgroups of $G$ with $N$ normal in $G$, let $NK = \{nk \mid n \in N, k \in K\}$.
}
    \textbf{(a).}
    Since $N$ is normal in $G$, then $N$ is also normal in $NK$ since $k \in K$ also means $k \in G$.

    \textbf{(b).}
    Given $a,b \in K$,
    \begin{align*}
        f(ab) &= Nab\\
        &= NaNb\\
        &= f(a)f(b).
    \end{align*}
    Thus $f$ is a homomorphism.
    Given $n \in N$, $k \in K$, then $Nnk = Nk = f(k)$, thus $f$ is surjective.
    We define
    \[
        \ker(f) = \{x \in K : f(x) = N\}.
    \]
    But note that $Nx = N$ means that also $x \in N$.
    Thus $K \cap N$ is our kernel, so $f$ is a surjective homomorphism with the restrictions as stated.

    \textbf{(c).}
    By \textbf{(b)}, the First Isomorphism Theorem gives
    \[
        K/(N \cap K) \cong NK/N.
    \]
\end{problem}

\begin{problem}[Exercise 18]{
    Suppose $N \trianglelefteq S_n$ such that $\ord(N)=2$.
}
    Then let $N=\{e,a\}$.
    Thus $ea = ae$ and $aa=e$.
    Hence $N \leq Z(S_n)$.
    By the lemma provided above, since $N \leq Z(S_n)$, we have $\ord(N) \leq \ord(Z(S_n))$.
    Hence a contradiction arises, since $2 \nleq 1$.
    Thus there is no subgroup of order $2$ in $S_n$ for $n > 3$.
\end{problem}

\begin{problem}[Exercise 19]{
    Given $N \leq S_n$ and $\sigma\tau = (1)$ for all $\sigma,\tau \neq (1)$ in $N$, determine $N$.
}
    Let $N \neq (1)$ and $\sigma \neq (1) \in N$.
    Then $\sigma\sigma = (1)$, which means $\sigma$ has order $2$.
    Thus for any other $\tau \neq (1)$ in $N$, we must have $\tau = \sigma$.
    Hence $N = (1)$ or $N$ is cyclic of order $2$.
\end{problem}

\begin{problem}[Exercise 20]{
    Given $N \trianglelefteq S_n$ and $N \cap A_n = A_n$, determine $N$.
}
    Then $A_n \subseteq N$ since $N \cap A_n = A_n$.
    Since $\ord(A_n) = n!/2$, and the order of $S_n$ is $n!$, $N$ can only have order $n!/2$ or $n!$.
    Hence $N$ is either $A_n$ or $S_n$.
\end{problem}